% Source-aligned English for printed p.206.
\clearpage
\SourcePage{206}

\LayerTarget{tohoku:proof:transformation-first-spectral-sequence:p1}
\noindent Theorem 5.3.1 we have
$\mathbf H^n(G,A_U)(y)=H^n(G,A_U(y))$ for $y\in Y_0$. By Proposition
5.2.3, the homomorphism
$H^n(Y,(A^G)_V)\to H^n(X;G,A_U)$ is therefore an isomorphism for
every $n$. It is just as easy to see that $H^n(X;G,A_{X_0})$ identifies with
$H^n(X_0;G,A)$; moreover, we know that
$H^n(Y,(A^G)_{Y_0})=H^n(Y_0,A^G)$. From the preceding diagram we then obtain,
by an evident ``diagram chase'', a sequence of homomorphisms

\LayerTarget{tohoku:sequence:5.4.2}
\[
\tag{5.4.2}
\begin{gathered}
{}\to H^{n-1}(X,G,A)
 \xrightarrow{\beta^{n-1}}
 H^{n-1}(X_0;G,A)/\operatorname{Im}H^{n-1}(Y_0,A^G)
 \xrightarrow{\partial} H^n(Y,A^G)\\[-0.2em]
 \xrightarrow{\alpha^n}H^n(X;G,A)
 \xrightarrow{\beta^n}
 H^n(X_0;G,A)/\operatorname{Im}H^n(Y_0,A^G)
 \to
\end{gathered}
\]
which is a complex (the product of two consecutive homomorphisms is zero) and
is exact everywhere except possibly at the terms $H^n(Y,A^G)$. The defect of
exactness $\operatorname{Ker}\alpha^n/\operatorname{Im}\partial$ at such a
term is canonically isomorphic to the kernel of the homomorphism
$H^n(Y_0,A^G)\to H^n(X_0;G,A)$. We thus obtain:

\LayerTarget{tohoku:proposition:5.4.1}
\textsc{Proposition 5.4.1.} \textit{The sequence of homomorphisms
\textup{(5.4.2)} defines on $H^n(Y,A^G)$ a composition series of length \textup{3} whose
successive quotients are respectively isomorphic to}
$\operatorname{Coker}\beta^{n-1}=H^{n-1}(X_0;G,A)/(\operatorname{Im}H^{n-1}(X;G,A)+\operatorname{Im}H^{n-1}(Y_0,A^G))$,
\textit{to}
$\operatorname{Ker}\bigl(H^n(Y_0,A^G)\to H^n(X_0;G,A)\bigr)$,
\textit{and to} $\operatorname{Ker}\beta^n\subset H^n(X;G,A)$.

\LayerTarget{tohoku:comparison:first-and-second-spectral-sequences-5.4}
Thus the first spectral sequence gives us a means of reducing the
determination of $H^*(Y,A^G)$ to good knowledge of $H^*(X,G,A)$,
$H^*(X_0,G,A)$, $H^*(Y_0,A^G)$, and the natural homomorphisms
$H^*(X;G,A)\to H^*(X_0;G,A)\leftarrow H^*(Y_0,A^G)$.
The determination of the first two groups can be attempted using the
\textit{second} spectral sequence, whereas the calculation of $H^*(Y_0,A^G)$
will often be possible if $Y_0$ can be chosen sufficiently small. A
particularly simple and important case is given by the following

\LayerTarget{tohoku:corollary:5.4.1}
\textsc{Corollary.} \textit{Suppose that $G$ acts trivially on $X_0$
(so $G$ is finite if $X_0$ is nonempty), that $A$ is a sheaf of vector spaces over a field $k$,
and that $G$ acts trivially on $A\mathbin{|}X_0$. Then the sequence \textup{(5.4.2)} is
exact, and $H^n(X_0,G,A)$ identifies canonically with}
$\displaystyle\bigoplus_{p+q=n}H^p(G,k)\otimes H^q(X_0,A)$.

\LayerTarget{tohoku:proof:corollary-5.4.1}
Indeed, the latter assertion is a special case of the last assertion of
Theorem 4.4.1. Consequently, the canonical homomorphism from
$H^n(Y_0,A^G)=H^n(X_0,A)$ into $H^n(X_0;G,A)$ is injective, so the sequence
(5.4.2) is exact.

\LayerTarget{tohoku:remark:smith-fixed-point-application:p1}
\textsc{Remarks.} Suppose that $G$ is a finite group of prime order.
Taking $X_0$ to be the set of points of $X$ fixed by $G$, and setting
$Y_0=f(X_0)$, the set $Y_0$ satisfies the condition stated at the beginning of
this section. If $A$ is then a sheaf of vector spaces arising from a sheaf on
$Y$ (i.e. $G$ acts trivially on $A\mathbin{|}X_0$), the hypotheses of
Corollary 1 hold. For example, one recovers very easily the theorem of P. A.
Smith concerning the case in which $X$ is a homology sphere mod. $p$ of
finite dimension ($p$ being the order of $G$): then $X_0$ is a homology sphere
mod. $p$. Moreover, this ``natural'' proof using general spectral sequences is
in fact very close to Borel's [2]; of course, the compactness hypothesis made

\TohokuContinuation{207}
